which is a functor on the category @@M0@@ (by (IV 2.12)). By (IV 3.2.3), this functor commutes with filtered inductive limits. Moreover, when @@M1@@ is a finite extension of @@M2@@, the projection @@M3@@ is a finite, flat morphism of finite presentation, and therefore defines a direct-image homomorphism @@M4@@. One immediately verifies that the condition on the norm homomorphisms considered in 7.18.3 does hold. Finally, if @@M5@@ is the coordinate algebra of an open subset @@M6@@ of the affine space @@M7@@, then @@M8@@ is identified with the map @@M9@@, which is surjective for @@M10@@ noetherian, by IX 1.6. Using the results proved above, one then obtains:
\begin{enumerate}[label=(\roman*)]
\item If @@M11@@ is a pure extension of @@M12@@, then @@M13@@ is injective, and is even bijective when @@M14@@ is noetherian.
\item If @@M15@@ is separably closed and @@M16@@ is a separable extension, then @@M17@@ is injective.
\item For every extension @@M18@@ of @@M19@@, the kernel of @@M20@@ is a torsion group.
\end{enumerate}
These results remain valid when one assumes only that @@M21@@ is noetherian, provided one restricts to extensions @@M22@@ of @@M23@@ that are of finite type (so that @@M24@@ is still noetherian). Moreover, essentially the same results hold for @@M25@@ in place of the functor @@M26@@, with the sole difference that in (i) one cannot assert bijectivity in general, for lack of the homotopy theorem for @@M27@@ in place of @@M28@@; to conclude bijectivity, one would have to assume that @@M29@@ is noetherian and regular, which implies that the groups @@M30@@ and @@M31@@ are the same for @@M32@@ and for @@M33@@, respectively.

